\FloatBarrier
\clearpage
\section{Conclusion}
\label{sec:conclusion}
The next research step is often determined by evidence just obtained. EAR implements this mode of progress in three workflows. Idea discovery reassesses candidates against close prior work. The mathematical inner loop revises research around gaps identified in review, while the outer loop tests execution strategies through complete outcomes. Rebuttal finds evidence relevant to reviewer concerns, retains useful drafts, and continues filling gaps. Research materials, decision reasons, and version changes are saved along the way, providing the basis for further work.

Historical records offer initial evidence for these mechanisms. The fixed manuscript subset receives better model assessments after revision. In its historical batch, the G3 selected strategy achieves more passes and fewer research invocations per pass. Rating-change prediction shows both binary discrimination and a weakness in detecting falls. Each result has a specific task, denominator, and evaluation condition. The cost of using an existing strategy also needs to be distinguished from the cost of continuing to search for strategies.

\begin{insight}
\textbf{Keep research moving; keep people on the key decisions}\par
Researchers set the direction and standards and check key results. EAR continues retrieval, investigation, and revision within those constraints, returning new work and the evidence needed to decide what follows.
\end{insight}

The ultimate test is the human time required to produce trustworthy research. Future evaluation should use matched tasks and consistent quality checks, recording time spent on reading, coordination, checking, repairs, failed-route rework, and final review. It should also account for full machine costs and calendar duration. These measurements are needed to establish how far EAR's design objective is achieved in practical use.

\Sub{The open-source project and public materials}
The \href{https://github.com/liyingze07-svg/EffectiveAutoResearch-EAR-}{EAR open-source repository} provides idea discovery, autonomous mathematical research, and rebuttal modules. Aggregate data, the evaluation protocol, and vector figures accompany this report to support statistical and interpretive checks; Section~\ref{sec:rebuttal-results} describes their scope and limitations. Readers can try the relevant workflows, check the reported results, and contribute further improvements.\src{1,10}
